\documentclass[main]{subfiles}


\title[Advanced Hyperparameter Optimization]{Advanced Hyperparameter Optimization}
\subtitle{The Big Picture}
\author[Katharina Eggensperger]{Katharina Eggensperger}
\date{}

%\institute{}
%\date{}

\begin{document}
	
	\maketitle
	

%----------------------------------------------------------------------
%----------------------------------------------------------------------
\begin{frame}[c]{Is BayesOpt the ideal method for HPO?}

(Vanilla) Bayesian Optimization is

\begin{itemize}
    \item theoretically grounded, well-established, and widely studied,
    \item built for sequential decision-making problems in which evaluations are expensive and the objective function is unknown,
    \item based on minimal assumption\footnote{theoretical guarantees only require that the true function lies in the RKHS induced by the kernel} which makes it widely applicable.
\end{itemize}

$\leadsto$ \textbf{Yes,} because it is more sample-efficient than random search, more adaptive than grid search, and more systematic than manual search. \\
$\leadsto$ \textbf{But}, it might still be too inefficient for practical HPO tasks.

\end{frame}
%-----------------------------------------------------------------------
%----------------------------------------------------------------------
\begin{frame}[c]{What do we know about HPO?}

Some observations:

\begin{itemize}
    \item we often search the same space on different dataset \\ $\leadsto$ similar tasks might have similar objective functions.
    \item prior information is available from previous experiments or the domain \\ $\leadsto$ we do not need to start from scratch.
    \item evaluating ML pipelines is iterative \\ $\leadsto$ we can get low-cost approximations of the true, costly function.
\end{itemize}

$\leadsto$ HPO has a lot of structure; treating it as generic black-box optimization is unnecessarily general and may be inefficient

\end{frame}
%-----------------------------------------------------------------------
%----------------------------------------------------------------------
\begin{frame}{Summary} 

Most important insights from this video:

\begin{enumerate}
    \item Bayesian optimization can be improved by exploiting the structure of the optimization problems.
    \item HPO offers task similarity, prior information, and intermediate results.
\end{enumerate}

$\leadsto$ We derive HPO-specific assumptions that increase the efficiency of Bayesian optimization for HPO tasks.

\end{frame}

\end{document}
